\documentclass[aspectratio=169,10pt]{beamer}

\usetheme{Madrid}
\usecolortheme{default}
\usefonttheme{professionalfonts}

\usepackage[utf8]{inputenc}
\usepackage[T1]{fontenc}
\usepackage[english]{babel}
\usepackage{amsmath,amssymb}
\usepackage{booktabs}
\usepackage{array}
\usepackage{graphicx}
\graphicspath{{./}{images/}{docs/images/}}
\setkeys{Gin}{draft=false}

\definecolor{GameGreen}{RGB}{0,105,92}
\definecolor{GameGold}{RGB}{225,151,48}
\definecolor{GameDark}{RGB}{33,43,54}
\definecolor{GameLight}{RGB}{243,247,246}

\setbeamercolor{structure}{fg=GameGreen}
\setbeamercolor{frametitle}{fg=white,bg=GameGreen}
\setbeamercolor{title}{fg=GameGreen}
\setbeamercolor{block title}{fg=white,bg=GameGreen}
\setbeamercolor{block body}{bg=GameLight,fg=GameDark}
\setbeamercolor{alertblock title}{fg=white,bg=GameGold}
\setbeamercolor{alertblock body}{bg=GameLight,fg=GameDark}
\setbeamercolor{author in head/foot}{fg=white,bg=GameGreen}
\setbeamercolor{title in head/foot}{fg=white,bg=GameDark}
\setbeamercolor{date in head/foot}{fg=white,bg=GameGreen}
\setbeamertemplate{navigation symbols}{}
\setbeamertemplate{itemize items}[triangle]
\setbeamertemplate{footline}{%
  \leavevmode\hbox{%
  \begin{beamercolorbox}[wd=.25\paperwidth,ht=2.5ex,dp=1ex,center]{author in head/foot}%
    Group 11
  \end{beamercolorbox}%
  \begin{beamercolorbox}[wd=.60\paperwidth,ht=2.5ex,dp=1ex,leftskip=.8em]{title in head/foot}%
    Horror Investigation Game
  \end{beamercolorbox}%
  \begin{beamercolorbox}[wd=.15\paperwidth,ht=2.5ex,dp=1ex,center]{date in head/foot}%
    \insertframenumber/\inserttotalframenumber
  \end{beamercolorbox}}%
  \vskip0pt%
}
\setbeamertemplate{title page}{%
  \vbox{}
  \vfill
  \begin{beamercolorbox}[wd=.96\paperwidth,sep=1.05em,center,rounded=true,shadow=true]{frametitle}
    {\usebeamerfont{title}\color{white}\inserttitle\par}
    \vspace{0.25cm}
    {\usebeamerfont{subtitle}\color{white}\insertsubtitle\par}
    \vspace{0.55cm}
    {\small\color{white}\insertinstitute\par}
    \vspace{0.20cm}
    {\small\color{white}\insertdate\par}
  \end{beamercolorbox}
  \vspace{1.0cm}
  \begin{center}
    \usebeamerfont{author}\insertauthor
  \end{center}
  \vfill
}

\title{Horror Investigation Game}
\subtitle{Technical Design and Prototype Implementation}
\author{\textbf{Group 11}\\[0.2cm]\small
Vo Minh Dang -- 23125002\\
Truong Gia Bao -- 23125052\\
Nguyen Nhat Khoa -- 23125085}
\institute{Advanced Program (APCS) \\ 3D Visualization and Game Development -- 23TT}
\date{Final Project Seminar S2: Tech -- 1 August 2026}

\begin{document}

\begin{frame}
  \titlepage
\end{frame}

\begin{frame}{Technical Goal: Build One Reliable Gameplay Foundation}
  \begin{block}{Engineering objective}
    The project first establishes a stable third-person gameplay framework in the dormitory vertical slice. Player, interaction, UI, inventory, objectives, and saving are designed to be reusable when later maps are integrated.
  \end{block}

  \vspace{0.35cm}
  \begin{columns}[T,onlytextwidth]
    \column{.48\textwidth}
    \begin{block}{Reusable systems}
      Player movement, camera, input, interaction, inventory, HUD, objective events, and save data.
    \end{block}
    \column{.48\textwidth}
    \begin{block}{Map-specific content}
      Scene geometry, puzzles, NPC placement, visual atmosphere, and the objectives that connect each chapter.
    \end{block}
  \end{columns}
\end{frame}

\section{Unity Foundation}
\begin{frame}{Unity Environment and Project Organization}
  \begin{columns}[T,onlytextwidth]
    \column{.48\textwidth}
    \begin{block}{Unity stack}
      \begin{itemize}
        \item Unity 6000.5.0f1 with Universal Render Pipeline 17.5.0.
        \item Input System 1.19.0 for keyboard, mouse, and gamepad actions.
        \item AI Navigation 2.0.13, Timeline 1.8.12, UGUI, and Unity Test Framework.
      \end{itemize}
    \end{block}
    \column{.48\textwidth}
    \begin{block}{Asset organization}
      \begin{itemize}
        \item \texttt{Assets/Chapter1}: runtime scripts, scene, UI, input, save, and tests.
        \item \texttt{Assets/Chapter2}: police station scene and 3D model.
        \item \texttt{Assets/Keypad}: reusable keypad and door assets.
      \end{itemize}
    \end{block}
  \end{columns}
\end{frame}

\begin{frame}{Runtime Architecture Separates Systems from Map Content}
  \begin{columns}[T,onlytextwidth]
    \column{.31\textwidth}
    \begin{block}{Player Layer}
      Input reader, motor, visual controller, stamina, input lock, camera target, and third-person camera rig.
    \end{block}
    \column{.31\textwidth}
    \begin{block}{Gameplay Layer}
      Interactable contract, interaction controller, pickups, inventory, doors, keypad, NPC logic, and objective markers.
    \end{block}
    \column{.31\textwidth}
    \begin{block}{Progression Layer}
      Chapter manager, event bus, chapter steps, HUD notifications, and JSON save service.
    \end{block}
  \end{columns}

  \vspace{0.35cm}
  \begin{alertblock}{Design benefit}
    A new scene can reuse the runtime layers while only replacing scene objects, chapter objectives, and encounter design.
  \end{alertblock}
\end{frame}

\section{Implemented Systems}
\begin{frame}{Player Controller and Camera}
  \begin{columns}[T,onlytextwidth]
    \column{.48\textwidth}
    \begin{block}{Character control}
      \begin{itemize}
        \item \texttt{CharacterController}-based third-person movement.
        \item Sprinting consumes stamina; sprint is unavailable until enough stamina recovers.
        \item Crouch supports low-ceiling checks and context-based input locks.
      \end{itemize}
    \end{block}
    \column{.48\textwidth}
    \begin{block}{Camera control}
      \begin{itemize}
        \item Orbit camera follows a dedicated target in \texttt{LateUpdate}.
        \item Camera collision uses a sphere cast against the environment.
        \item Look input can be disabled during dialogue, puzzles, or pause states.
      \end{itemize}
    \end{block}
  \end{columns}
\end{frame}

\begin{frame}{Interaction, Inventory, and Door Systems}
  \begin{columns}[T,onlytextwidth]
    \column{.48\textwidth}
    \begin{block}{Interaction pipeline}
      \begin{itemize}
        \item Camera-forward sphere cast detects targets on the \texttt{Interactable} layer.
        \item \texttt{IChapter1Interactable} provides a common contract for world objects.
        \item Prompts, highlights, one-shot behavior, and step gating are handled consistently.
      \end{itemize}
    \end{block}
    \column{.48\textwidth}
    \begin{block}{Items and access control}
      \begin{itemize}
        \item Pickups update \texttt{PlayerInventory} and persisted unique item IDs.
        \item Flashlight, fuse, throwable can, and inspectable objects are supported.
        \item Room-door keypad logic controls password-gated access from the outside.
      \end{itemize}
    \end{block}
  \end{columns}
\end{frame}

\begin{frame}{Objectives, HUD, and Save Data}
  \begin{columns}[T,onlytextwidth]
    \column{.48\textwidth}
    \begin{block}{Event-driven presentation}
      \begin{itemize}
        \item \texttt{Chapter1EventBus} publishes inventory and progression changes.
        \item HUD modules cover objectives, interaction prompts, inventory, stamina, notifications, and debug information.
        \item Scene markers provide named objectives without blocking player movement.
      \end{itemize}
    \end{block}
    \column{.48\textwidth}
    \begin{block}{Persistent progression}
      \begin{itemize}
        \item \texttt{Chapter1SaveData} stores only serializable gameplay state.
        \item \texttt{JsonChapter1SaveService} saves JSON under \texttt{Application.persistentDataPath}.
        \item World pickups, inventory changes, and chapter state remain synchronized.
      \end{itemize}
    \end{block}
  \end{columns}
\end{frame}

\begin{frame}{Henry Pursuit Prototype: Animation and Gameplay Integration}
  \begin{columns}[T,onlytextwidth]
    \column{.48\textwidth}
    \centering
    \includegraphics[height=.40\textheight,keepaspectratio]{HenryRunPose_Quarter.png}

    \vspace{0.10cm}
    \small\textbf{Front view:} corrected run pose on Henry's production model.

    \column{.48\textwidth}
    \centering
    \includegraphics[height=.40\textheight,keepaspectratio]{HenryRunPose_ThreeQuarter.png}

    \vspace{0.10cm}
    \small\textbf{Three-quarter view:} stride and body-stability validation.
  \end{columns}

  \vspace{0.10cm}
  \begin{alertblock}{Runtime behavior}
    Henry uses idle and run animation controllers. After the theft interaction, the chase controller switches the NPC into the pursuit state and presents a game-over result when the player is caught.
  \end{alertblock}
\end{frame}

\section{Scene Integration and Quality}
\begin{frame}{Scene Readiness and Technical Scope}
  \renewcommand{\arraystretch}{1.25}
  \begin{tabular}{p{.20\linewidth} p{.37\linewidth} p{.31\linewidth}}
    \toprule
    \textbf{Chapter} & \textbf{Current technical state} & \textbf{Next integration target} \\
    \midrule
    Dormitory & Playable scene with player, interaction, inventory, HUD, save, keypad, objective markers, and NPC-related systems. & Stabilize the complete vertical slice and chapter exit. \\
    Police Station & Scene and 3D model are present. & Add chapter objectives, archive puzzle, stealth zones, and scene transition. \\
    Old Hospital & Dedicated chapter assets and scene are not yet present. & Build horror encounters, power restoration, checkpoints, and evidence flow. \\
    Killer's House & Dedicated chapter assets and scene are not yet present. & Build final stealth encounter, endings, and final transition. \\
    \bottomrule
  \end{tabular}
\end{frame}

\begin{frame}{Validation Strategy Focuses on Progression Blockers}
  \begin{columns}[T,onlytextwidth]
    \column{.48\textwidth}
    \begin{block}{Automated edit-mode checks}
      \begin{itemize}
        \item Required dormitory areas, spawns, and objective markers exist.
        \item Colliders are present where required and markers do not block movement.
        \item Room-door keypad, camera, and audio-listener configuration are valid.
        \item The scene contains no missing-script components.
      \end{itemize}
    \end{block}
    \column{.48\textwidth}
    \begin{block}{Manual play-mode checks}
      \begin{itemize}
        \item Movement, camera collision, sprint stamina, crouch, and input locks.
        \item Interaction focus, inventory feedback, keypad access, and save/load.
        \item NPC chase and failure flow, then full chapter progression to the exit.
      \end{itemize}
    \end{block}
  \end{columns}
\end{frame}

\section{Technical Delivery Plan}
\begin{frame}{Technical Roadmap Before Final Delivery}
  \begin{enumerate}
    \item Finish and verify the Chapter 1 vertical slice from opening objective to scene exit.
    \item Generalize shared player, interaction, UI, and save interfaces for cross-chapter reuse.
    \item Integrate the police-station objective chain, puzzle gates, and stealth requirements.
    \item Build hospital and killer-house scenes with checkpoint, threat, and ending systems.
    \item Add playable scenes to Build Settings and perform full-path regression testing.
  \end{enumerate}

  \vspace{0.2cm}
  \begin{alertblock}{Delivery check}
    Verify the final build, project reopening, main progression path, and save/load behavior before packaging \texttt{Group11-S2.rar}.
  \end{alertblock}
\end{frame}

\begin{frame}{Technical Responsibilities}
  \begin{columns}[T,onlytextwidth]
    \column{.31\textwidth}
    \begin{block}{Vo Minh Dang}
      \begin{minipage}[t][.35\textheight][t]{\linewidth}
        \small
        \textbf{Gameplay \& Integration}
        \begin{itemize}
          \item Interaction, objectives, and NPC/AI systems.
          \item Event connections between gameplay and UI.
          \item System integration and bug fixing.
        \end{itemize}
      \end{minipage}
    \end{block}

    \column{.34\textwidth}
    \begin{alertblock}{Truong Gia Bao -- UI/HUD}
      \begin{minipage}[t][.35\textheight][t]{\linewidth}
        \small
        \textbf{UI Implementation}
        \begin{itemize}
          \item Interaction prompts and notifications.
          \item Objective, inventory, and stamina HUD.
          \item UGUI layout, readability, and consistency.
        \end{itemize}
      \end{minipage}
    \end{alertblock}

    \column{.31\textwidth}
    \begin{block}{Nguyen Nhat Khoa}
      \begin{minipage}[t][.35\textheight][t]{\linewidth}
        \small
        \textbf{Content \& QA}
        \begin{itemize}
          \item Objective and notification content.
          \item Gameplay-flow and UI validation.
          \item Issue tracking and demo preparation.
        \end{itemize}
      \end{minipage}
    \end{block}
  \end{columns}

  \vspace{0.15cm}
  \small\centering
  \textbf{Technical workflow:} Bao implements the UI, Dang connects it to gameplay, and Khoa validates the user-facing flow.
\end{frame}

\begin{frame}{Conclusion: A Reusable Path from Prototype to Full Game}
  \begin{block}{Technical takeaway}
    The dormitory prototype proves the core gameplay stack: player control, interaction, inventory, UI, save data, access control, and NPC pursuit. The next phase is scene integration and validation so the same foundation can support the remaining chapters.
  \end{block}

  \vfill
  \begin{center}
    \Large\textcolor{GameGreen}{\textbf{Build the system once. Validate it thoroughly. Reuse it across every map.}}
  \end{center}
\end{frame}

\end{document}
